\documentclass{article}
% A tiny fixture paper for the second-brain test suite.
\usepackage{amsmath,amsthm}

% Shared notation, inlined via \input.
\newcommand{\mech}{\mathcal{M}}


\newcommand{\val}{v_i(t_i)}
\newcommand{\typespace}[1]{T_{#1}}
\providecommand{\eps}{\varepsilon}
\DeclareMathOperator*{\argmax}{arg\,max}

\newtheorem{theorem}{Theorem}
\newtheorem{lemma}[theorem]{Lemma}
\newtheorem{definition}{Definition}
\newtheorem{assumption}{Assumption}

\title{Allocation with Monotone Valuations: A Toy Study\thanks{Fixture only.}}
\author{A. Author \and B. Author}
\date{January 2026}

\begin{document}
\maketitle

\begin{abstract}
We study a toy allocation problem with monotone valuations $\val$ and show
that the direct mechanism $\mech$ is optimal. % trailing comment to strip
\end{abstract}

\section{Introduction}
Allocation problems with private types are classical. Our contribution is a
self-contained fixture: every result below is intentionally trivial, but the
notation ($\val$, $\mech$, $\typespace{i}$) exercises macro handling, and the
cross-references exercise dependency mapping. 50\% of this sentence tests
escaped percent signs.

\section{Model}

\begin{definition}[Mechanism]\label{def:mechanism}
A \emph{mechanism} $\mech = (x, p)$ consists of an allocation rule
$x : \typespace{1} \times \dots \times \typespace{n} \to [0,1]^n$ and a
payment rule $p : \typespace{1} \times \dots \times \typespace{n} \to \mathbb{R}^n$.
\end{definition}

\begin{definition}[Valuation]\label{def:valuation}
Agent $i$'s \emph{valuation} for the good is $\val$, a function of her type
$t_i \in \typespace{i}$.
\end{definition}

\begin{assumption}[Monotonicity]\label{ass:mono}
For every agent $i$, the valuation $\val$ is nondecreasing in $t_i$.
\end{assumption}

\section{Main Results}

\begin{theorem}[Optimality of the direct mechanism]\label{thm:main}
Under Assumption~\ref{ass:mono}, the direct mechanism $\mech^{*}$ with
$x^{*}(t) \in \argmax_{x} \sum_i \val$ maximizes welfare among all mechanisms
of Definition~\ref{def:mechanism}.
\end{theorem}

\begin{proof}
Fix a type profile $t$. By Definition 2, each $\val$ depends only on $t_i$, so
welfare is separable. By Assumption 1, pointwise maximization is monotone in
each coordinate, hence $x^{*}$ is well defined and feasible, and no mechanism
of Definition~\ref{def:mechanism} achieves higher welfare. Take $\eps \to 0$.
\end{proof}

\begin{lemma}[Payment normalization]\label{lem:aux}
Under Assumption~\ref{ass:mono}, payments can be normalized so that
$p_i(0, t_{-i}) = 0$ for all $i$ without changing the allocation of
Theorem~\ref{thm:main}.
\end{lemma}

\begin{proof}[Proof of Lemma~\ref{lem:aux}]
Subtract $p_i(0, t_{-i})$ from agent $i$'s payment; the allocation rule of
Theorem 1 is unchanged.
\end{proof}

\section{Conclusion}
A fixture is not a contribution.

\begin{thebibliography}{9}
\bibitem{myerson1981} R. Myerson. Optimal auction design. \emph{Mathematics of
Operations Research}, 6(1):58--73, 1981.
\bibitem{nisan2006} N. Nisan and I. Segal. The communication requirements of
efficient allocations. \emph{Journal of Economic Theory}, 129(1):192--224, 2006.
\end{thebibliography}

\end{document}
